\begin{tikzpicture}[
  x=1cm,y=1cm,font=\small,
  pgm variable/.style={circle,draw=BookInk,fill=BookPaper,minimum size=7mm,inner sep=1pt},
  pgm relation/.style={draw=BookTeal,line width=0.8pt}
]
  \node[text=BookInk] at (1.2,2.1) {目标三角图};
  \node[pgm variable] (a1) at (1.2,1.35) {$Z_1$};
  \node[pgm variable] (a2) at (0.35,0) {$Z_2$};
  \node[pgm variable] (a3) at (2.05,0) {$Z_3$};
  \draw[pgm relation] (a1)--(a2)--(a3)--(a1);
  \node[align=center,text=BookMuted] at (1.2,-0.85)
    {$\Pr(Z_i\ne Z_j)=8/13$\\可实现的相关性};

  \draw[BookRule] (2.9,-1.3)--(2.9,2.4);
  \node[text=BookInk] at (4.55,2.1) {完全平均场};
  \node[pgm variable] (b1) at (4.55,1.35) {$Z_1$};
  \node[pgm variable] (b2) at (3.7,0) {$Z_2$};
  \node[pgm variable] (b3) at (5.4,0) {$Z_3$};
  \node[text=BookBlue] at (4.55,0.65) {$q_1q_2q_3$};
  \node[align=center,text=BookMuted] at (4.55,-0.85)
    {$q(Z_i\ne Z_j)=1/2$\\删去全部依赖};

  \draw[BookRule] (6.2,-1.3)--(6.2,2.4);
  \node[text=BookInk] at (7.85,2.1) {局部LP分数解};
  \node[pgm variable] (c1) at (7.85,1.35) {$Z_1$};
  \node[pgm variable] (c2) at (7,0) {$Z_2$};
  \node[pgm variable] (c3) at (8.7,0) {$Z_3$};
  \draw[BookAmber,densely dashed,line width=0.8pt] (c1)--(c2)--(c3)--(c1);
  \node[text=BookAmber] at (7.85,0.5) {$\ne$};
  \node[align=center,text=BookMuted] at (7.85,-0.85)
    {$\mu(Z_i\ne Z_j)=1$\\三条边无法同时满足};
\end{tikzpicture}
